\chapter{Agent Architectures and Orchestration: ReAct, Planning, and Multi-Agent Systems}
\label{ch:agent_architectures}

\begin{chaptergoalsbox}
\textbf{What You Will Learn:}
\begin{itemize}[leftmargin=15pt]
  \item Core cognitive agent loops: ReAct (Reason + Act), Plan-and-Solve, and Reflexion verbal reinforcement.
  \item Strict tool-calling mechanics: JSON Schema definitions, OpenAI/Anthropic wire protocols, and Pydantic argument validation.
  \item Deterministic execution runtimes: Fault handling, dynamic exception reflection, rate-limiting tripwires, and loop detectors.
  \item Multi-agent topologies: Centralized Supervisor-Worker hierarchies, debate/consensus protocols, and state isolation barriers.
  \item Production safety engineering: Guardrails, privilege sandboxing, and token budget ceilings.
\end{itemize}

\textbf{What You Should Already Understand:}
\begin{itemize}[leftmargin=15pt]
  \item JSON Schema specification standard (Draft 7/2020-12).
  \item Modern Python asynchronous programming (\texttt{asyncio}, typing, error handling).
  \item REST APIs and HTTP status codes.
\end{itemize}

\textbf{Why This Topic Exists:}
A single LLM prompt cannot perform real work in external environments. Autonomous agents must interact with operating systems, databases, and third-party APIs. Without robust architectural scaffolding, agents fall into infinite loops, hallucinate non-existent tool arguments, leak customer data across boundaries, or crash catastrophically upon receiving an unhandled 500 error from an API.

\textbf{Motivating Practical Problem:}
A DevOps remediation agent is tasked with diagnosing a failing Kubernetes cluster. The agent must inspect pod crash logs, query Prometheus metrics, identify that an out-of-memory (OOM) error occurred, calculate a safe memory limit increase, update the Helm values repository via a Git pull request, and verify pod health. If the Prometheus query times out, the agent must not crash; it must reflect on the network timeout, switch to an alternative logging endpoint, and complete the remediation safely.
\end{chaptergoalsbox}

\section{Cognitive Agent Paradigms}

An AI agent is an autonomous entity that perceives its environment through sensors (text prompts, API responses, tool outputs) and acts upon that environment through effectors (tool executions, code generation, database writes) to achieve a specified goal.

\subsection{The ReAct Pattern: Synergizing Reasoning and Acting}
Prior to ReAct (Yao et al., 2022), models either reasoned in isolation (Chain-of-Thought) or acted greedily without reasoning (direct tool dispatch). 
ReAct interleaves internal reasoning with external actions:
\begin{equation}
\text{Thought}_t \longrightarrow \text{Action}_t \longrightarrow \text{Observation}_t \longrightarrow \text{Thought}_{t+1} \longrightarrow \dots
\end{equation}

\begin{figure}[htbp]
\centering
\begin{tikzpicture}[
  node/.style={draw=headingblue, fill=codebg, rounded corners=2pt, font=\sffamily\scriptsize\bfseries, minimum width=3.2cm, minimum height=1.0cm, align=center},
  arrow/.style={-{Stealth[scale=0.9]}, thick, color=linkblue}
]

\node[node, fill=intuitionBg] (thought) at (0, 2) {1. Thought\\Deconstruct task \& plan next step};
\node[node, fill=exampleBg] (action) at (5, 2) {2. Action\\Dispatch tool call with parameters};
\node[node, fill=warningBg] (env) at (5, -0.2) {3. Environment\\Execute API / DB / Shell command};
\node[node, fill=gray!10] (obs) at (0, -0.2) {4. Observation\\Parse output into conversation context};

\draw[arrow] (thought) -- (action);
\draw[arrow] (action) -- (env);
\draw[arrow] (env) -- (obs);
\draw[arrow] (obs) -- (thought);

\end{tikzpicture}
\caption{The Cyclic ReAct (Reason + Act) Execution Loop.}
\label{fig:react_loop}
\end{figure}

The reasoning trace (\textbf{Thought}) allows the model to induce, track, and update action plans as well as handle exceptions. The external feedback (\textbf{Observation}) grounds the model's reasoning in physical reality, preventing hallucinations.

\subsection{Plan-and-Solve: Mitigating Reasoning Drift}
While ReAct is flexible, it suffers from "myopic drift": on complex 20-step tasks, the model frequently forgets the original overarching objective and gets sidetracked by secondary issues.
Plan-and-Solve (Wang et al., 2023) decouples execution into two distinct phases:
\begin{enumerate}[leftmargin=18pt]
  \item \textbf{Planner Model}: Decomposes the complex user goal into an explicit Directed Acyclic Graph (DAG) of discrete sub-tasks.
  \item \textbf{Executor Model}: Iteratively executes each sub-task in sequence, feeding intermediate results to a Coordinator that evaluates whether the plan requires dynamic replanning.
\end{enumerate}

\subsection{Reflexion: Episodic Memory and Self-Correction}
Reflexion (Shinn et al., 2023) equips agents with dynamic memory and self-reflective capabilities:

\begin{figure}[htbp]
\centering
\begin{tikzpicture}[
  box/.style={draw=headingblue, fill=codebg, rounded corners=2pt, font=\sffamily\scriptsize, minimum width=3.0cm, minimum height=1.0cm, align=center},
  arrow/.style={-{Stealth[scale=0.8]}, thick, color=headingblue}
]

\node[box] (actor) at (0, 1.5) {Actor\\Generates trajectory};
\node[box, fill=exampleBg] (eval) at (4.0, 1.5) {Evaluator\\Scores performance};
\node[box, fill=warningBg] (self) at (8.0, 1.5) {Self-Reflection\\Identifies failure cause};
\node[box, fill=intuitionBg] (mem) at (4.0, -0.5) {Episodic Memory Buffer\\Stores lessons learned};

\draw[arrow] (actor) -- (eval);
\draw[arrow] (eval) -- (self);
\draw[arrow] (self) |- (mem);
\draw[arrow] (mem) -| (actor);

\end{tikzpicture}
\caption{Reflexion Architecture: Failed execution traces trigger verbal self-reflection stored in episodic memory to guide subsequent trials.}
\label{fig:reflexion_arch}
\end{figure}

When a trajectory fails (e.g., a unit test fails or a compiler emits an error), the agent does not immediately re-run the same prompt. Instead, a dedicated Reflection prompt produces a verbal post-mortem:
\begin{center}
\emph{"In my previous attempt, I passed an integer string instead of a UNIX timestamp to the API. In the next trial, I will explicitly convert the ISO-8601 string using \texttt{datetime.timestamp()}."}
\end{center}
This reflection is appended to the agent's short-term episodic memory buffer, dramatically reducing repeated errors across complex tasks.

\section{Tool Calling Protocols and JSON Schema Validation}

Modern frontier models are trained natively to emit structured tool invocation calls using JSON Schema specifications.

\subsection{The Wire Protocol}
When tools are provided to an LLM, the engine serializes them into the request payload:
\begin{lstlisting}[language=Python, caption={Declaring Tool Schemas with Pydantic}]
from pydantic import BaseModel, Field
from typing import Literal

class RestartPodInput(BaseModel):
    namespace: str = Field(
        ..., 
        description="The Kubernetes namespace containing the target pod, e.g., 'production'."
    )
    pod_name: str = Field(
        ..., 
        description="The exact alphanumeric name of the pod to restart."
    )
    grace_period_seconds: int = Field(
        default=30, 
        ge=0, 
        le=300, 
        description="Grace period in seconds before SIGKILL is dispatched."
    )
    force: bool = Field(
        default=False, 
        description="Set to true only if standard deletion fails."
    )

# The generated JSON Schema conforms to standard Draft 7
tool_definition = {
    "type": "function",
    "function": {
        "name": "restart_kubernetes_pod",
        "description": "Restarts a misbehaving pod by dispatching a Kubernetes delete command.",
        "parameters": RestartPodInput.model_json_schema()
    }
}
\end{lstlisting}

\subsection{Handling Wire Output and Type Ingestion}
Frontier LLMs return structured \texttt{tool\_calls} objects. A production agent must never trust the raw output without schema validation:
\begin{lstlisting}[language=Python, caption={Safe Tool Dispatch and Validation Loop}]
import json
from pydantic import ValidationError

async def dispatch_tool_call(tool_call, registered_tools: dict) -> dict:
    func_name = tool_call.name
    raw_args = tool_call.arguments
    
    if func_name not in registered_tools:
        return {
            "is_error": True,
            "error": f"Tool '{func_name}' does not exist. Available tools: {list(registered_tools.keys())}"
        }
        
    tool_meta = registered_tools[func_name]
    schema_model = tool_meta["schema"]
    callable_func = tool_meta["handler"]
    
    # 1. Parse JSON string
    try:
        args_dict = json.loads(raw_args) if isinstance(raw_args, str) else raw_args
    except json.JSONDecodeError as e:
        return {
            "is_error": True,
            "error": f"Invalid JSON syntax in arguments: {str(e)}. Please correct your JSON output."
        }
        
    # 2. Validate against Pydantic schema
    try:
        validated_args = schema_model(**args_dict)
    except ValidationError as ve:
        return {
            "is_error": True,
            "error": f"Parameter validation failed: {ve.json()}. Adjust parameters to match schema."
        }
        
    # 3. Execute tool within safety boundaries
    try:
        result = await callable_func(validated_args)
        return {"is_error": False, "result": result}
    except Exception as exc:
        # Capture runtime exceptions to feed back to the LLM
        return {
            "is_error": True,
            "error": f"Execution error in '{func_name}': {type(exc).__name__}: {str(exc)}"
        }
\end{lstlisting}

\begin{underthehoodbox}[Error Feedback as Model Guidance]
Notice that when \texttt{ValidationError} or execution exceptions occur, we do \textbf{not} crash the agent process. Instead, we serialize the error into an \texttt{Observation} message tagged with the \texttt{tool\_call\_id}.
Modern frontier LLMs are remarkably adept at reading Python tracebacks and schema validation errors; upon receiving the error message, the model inspects the missing field, adjusts its parameters, and re-invokes the tool correctly on the next turn.
\end{underthehoodbox}

\section{Deterministic Tool Execution Runtimes}

An autonomous agent operating without safety guards is a catastrophic liability. Production runtimes enforce three non-negotiable boundaries:

\begin{enumerate}[leftmargin=18pt]
  \item \textbf{Infinite Loop Detection}: Models frequently enter circular reasoning (e.g., Tool A fails $\to$ calls Tool B $\to$ fails $\to$ calls Tool A). Track a rolling hash of the last 4 tool calls and arguments; if a duplicate sequence is detected, inject a hard intervention prompt.
  \item \textbf{Max-Iteration and Timeout Tripwires}: Set strict maximum step counts (e.g., $N \le 12$) and wall-clock execution timeouts (e.g., 60 seconds).
  \item \textbf{Token Budget Ceilings}: Running an autonomous loop that writes 5,000-token tool outputs into the conversation history can exhaust context windows or incur hundreds of dollars per session. Truncate tool observations to a strict token ceiling (e.g., max 1,500 tokens).
\end{enumerate}

\begin{lstlisting}[language=Python, caption={Production Autonomous Execution Loop with Tripwires}]
import asyncio
from typing import Sequence

class ProductionAgentLoop:
    def __init__(self, llm_client, tools: dict, max_steps: int = 10, max_tokens: int = 15000):
        self.llm = llm_client
        self.tools = tools
        self.max_steps = max_steps
        self.max_tokens = max_tokens

    async def run(self, user_prompt: str) -> str:
        messages = [{"role": "user", "content": user_prompt}]
        step_count = 0
        total_tokens_consumed = 0
        call_history = []
        
        while step_count < self.max_steps:
            step_count += 1
            
            # Step 1: Model inference
            response = await self.llm.invoke(messages, tools=list(self.tools.values()))
            total_tokens_consumed += response.usage.total_tokens
            
            if total_tokens_consumed > self.max_tokens:
                raise TimeoutError("Token budget exceeded safety limit. Aborting agent execution.")
                
            messages.append(response.message)
            
            # If the model produced a final answer without requesting tools, exit
            if not response.tool_calls:
                return response.content
                
            # Step 2: Process tool calls
            for tool_call in response.tool_calls:
                call_sig = (tool_call.name, tool_call.arguments)
                if call_history.count(call_sig) >= 2:
                    # Loop detected: inject warning directly into observation
                    tool_result = {
                        "is_error": True,
                        "error": "CRITICAL: You have invoked this exact tool with identical parameters twice. You must alter your strategy."
                    }
                else:
                    call_history.append(call_sig)
                    tool_result = await dispatch_tool_call(tool_call, self.tools)
                    
                # Truncate output to prevent prompt bloating
                str_res = str(tool_result["result"] if not tool_result["is_error"] else tool_result["error"])
                if len(str_res) > 2000:
                    str_res = str_res[:2000] + "... [TRUNCATED DUE TO SIZE]"
                    
                messages.append({
                    "role": "tool",
                    "tool_call_id": tool_call.id,
                    "content": str_res
                })
                
        return "Task could not be completed within maximum allowed reasoning steps."
\end{lstlisting}

\section{Multi-Agent Topologies}

When systems exceed the reasoning capacity of a single context window, multi-agent architectures partition responsibilities across specialized cooperating agents.

\subsection{Centralized Supervisor-Worker Architecture}
A top-level \textbf{Supervisor} agent acts as a dispatcher. It maintains the user-facing interaction, decomposes requests, routes work to specialized domain workers, and aggregates results.

\begin{figure}[htbp]
\centering
\begin{tikzpicture}[
  supervisor/.style={draw=headingblue, fill=codebg, rounded corners=3pt, font=\sffamily\small\bfseries, minimum width=4.5cm, minimum height=1.2cm, align=center},
  worker/.style={draw=exampleBlue, fill=exampleBg, rounded corners=2pt, font=\sffamily\scriptsize, minimum width=3.2cm, minimum height=1.0cm, align=center},
  arrow/.style={-{Stealth[scale=0.9]}, thick, color=headingblue}
]

\node[supervisor] (sup) at (0, 2.5) {Supervisor / Router Agent\\(Global State \& Planning)};

\node[worker] (w1) at (-4.0, 0) {Research Worker\\(Web / Doc Retrieval)};
\node[worker] (w2) at (0, 0) {SQL Worker\\(Internal DB Queries)};
\node[worker] (w3) at (4.0, 0) {Code Sandbox Worker\\(Python REPL / Execution)};

\draw[arrow] (sup) -- (w1) node[midway, left, font=\tiny\sffamily] {Task Dispatch};
\draw[arrow] (w1) -- (sup) node[midway, right, font=\tiny\sffamily] {Summary};
\draw[arrow] (sup) -- (w2);
\draw[arrow] (w2) -- (sup);
\draw[arrow] (sup) -- (w3);
\draw[arrow] (w3) -- (sup);

\end{tikzpicture}
\caption{Supervisor-Worker Multi-Agent Orchestration.}
\label{fig:supervisor_worker}
\end{figure}

\begin{underthehoodbox}[Context Isolation in Multi-Agent Systems]
The single greatest advantage of the Supervisor-Worker topology is \textbf{context window isolation}. 
If the SQL worker executes 15 iterative queries against PostgreSQL to discover the correct schema, its internal execution trace generates 8,000 tokens of noisy metadata.
Instead of polluting the Supervisor's prompt with all 15 queries, the SQL Worker executes in an isolated context and returns only a clean, 200-token synthesized summary back to the Supervisor. The Supervisor's context window remains compact, pristine, and focused on global reasoning.
\end{underthehoodbox}

\begin{realworldbox}[Multi-Agent Code Review \& Security Hardening]
A financial enterprise implemented a three-agent peer-review consensus network:
\begin{enumerate}
  \item \textbf{Coder Agent}: Drafts infrastructure-as-code (Terraform) scripts.
  \item \textbf{Security Auditor Agent}: Scans generated code against CIS benchmarks and OWASP guidelines, emitting critique points.
  \item \textbf{Compliance Judge}: Evaluates the debate between Coder and Auditor.
\end{enumerate}
In testing, the Coder agent initially used public S3 buckets for temporary logs. The Security Auditor flagged the finding with rule \texttt{CIS-AWS-2.1.4.1}, prompting the Coder to inject KMS encryption keys and block public ACLs before any pull request was submitted.
\end{realworldbox}

\begin{tryityourselfbox}[Implement a Safe Calculator Tool with Sandboxing]
Build a tool that evaluates mathematical expressions requested by an LLM:
\begin{enumerate}
  \item Define a Pydantic model \texttt{MathExpression} with input string \texttt{expression}.
  \item \textbf{Anti-Pattern}: Do \emph{not} use Python's built-in \texttt{eval()}, which allows arbitrary code execution (\texttt{\_\_import\_\_('os').system('rm -rf /')}).
  \item Parse the expression using Python's \texttt{ast} module (Abstract Syntax Tree) and evaluate only safe binary operators (\texttt{Add, Sub, Mult, Div, Pow}).
  \item Return a structured JSON response. Verify that malicious inputs fail safely.
\end{enumerate}
\end{tryityourselfbox}

\section{Chapter Summary \& Key Takeaways}
\begin{itemize}[leftmargin=18pt]
  \item \textbf{Cognitive Architecture}: ReAct interleaves Thought, Action, and Observation, preventing reasoning detachment. Plan-and-Solve decomposes complex goals to prevent goal drift, while Reflexion uses verbal post-mortems stored in episodic memory to self-heal.
  \item \textbf{Tool Protocol}: Tools must be declared using standard JSON Schema and validated via Pydantic. Never pass raw string arguments directly to internal APIs.
  \item \textbf{Fault Resilience}: When tool calls fail, feed the exception back to the LLM context as a tool observation. Modern LLMs inspect tracebacks and self-correct on subsequent turns.
  \item \textbf{Tripwires \& Guards}: Enforce hard constraints on max reasoning steps, token budgets, and loop detectors to prevent runaway costs and infinite recursion.
  \item \textbf{Multi-Agent Isolation}: Centralized Supervisor-Worker architectures maintain high performance by isolating domain workers in independent context windows, returning only condensed summaries to the orchestrator.
\end{itemize}

\section{Exercises}

\subsection*{Mechanics}
\textbf{Exercise 5.1}: Convert the following Python function into an exact OpenAI-compatible JSON Schema tool specification:
\begin{lstlisting}[language=Python]
def transfer_funds(source_account: str, destination_account: str, amount_cents: int, currency: str = "USD") -> bool:
    """Transfers funds between corporate ledger accounts."""
    pass
\end{lstlisting}

\textbf{Exercise 5.2}: An agent executes a ReAct loop. In step 1, it calls \texttt{search\_database("user\_id=12")}. The tool returns \texttt{Error: Table 'users' does not contain column 'user\_id'. Available columns: ['id', 'email', 'created\_at']}. Write the exact JSON structure of the tool message that should be appended to the conversation history to trigger self-correction.

\subsection*{Debugging}
\textbf{Exercise 5.3}: An engineer connects an autonomous agent to a Slack bot. Within 2 hours of deployment, the company's OpenAI bill surges by \$4,000. Analysis indicates that two agents engaged in an infinite conversation, repeatedly thanking each other and asking if further help was required. Diagnose the failure in termination criteria and write the code guard to prevent it.

\subsection*{Architecture}
\textbf{Exercise 5.4}: Design a high-reliability Customer Service Multi-Agent System for an airline. The system consists of: a Triage Supervisor, a Flight Booking Worker, a Baggage Tracking Worker, and a Compensation Claims Worker. Define the state-sharing protocol, tool boundaries, and human escalation gates.

\section{Detailed Solutions}

\subsection*{Solution 5.1}
The standard JSON Schema tool definition is:
\begin{lstlisting}[language=Python]
tool_schema = {
    "type": "function",
    "function": {
        "name": "transfer_funds",
        "description": "Transfers funds between corporate ledger accounts.",
        "parameters": {
            "type": "object",
            "properties": {
                "source_account": {
                    "type": "string",
                    "description": "The originating corporate account identifier."
                },
                "destination_account": {
                    "type": "string",
                    "description": "The destination corporate account identifier."
                },
                "amount_cents": {
                    "type": "integer",
                    "description": "Amount to transfer in integer cents (e.g. 1000 for $10.00)."
                },
                "currency": {
                    "type": "string",
                    "enum": ["USD", "EUR", "GBP"],
                    "default": "USD",
                    "description": "Three-letter ISO currency code."
                }
            },
            "required": ["source_account", "destination_account", "amount_cents"]
        }
    }
}
\end{lstlisting}

\subsection*{Solution 5.2}
The message must be formatted with role \texttt{tool} and match the initiating \texttt{tool\_call\_id}:
\begin{lstlisting}[language=Python]
tool_message = {
    "role": "tool",
    "tool_call_id": "call_abc12345xyz",
    "name": "search_database",
    "content": json.dumps({
        "status": "error",
        "error_type": "InvalidColumnError",
        "message": "Table 'users' does not contain column 'user_id'. Available columns: ['id', 'email', 'created_at']"
    })
}
\end{lstlisting}

\subsection*{Solution 5.3}
The system lacked:
\begin{enumerate}
  \item \textbf{Turn-Limit Counter}: No monotonic step counter bounded the conversation.
  \item \textbf{Explicit Termination Token}: The prompt failed to instruct the model to emit a terminal keyword (such as \texttt{[TERMINATE]}) when no actionable task remained.
\end{enumerate}
Fix: Enforce a maximum turn limit and check for termination tokens:
\begin{lstlisting}[language=Python]
if "[TASK_COMPLETE]" in response.content or step_count >= 6:
    return clean_response(response.content)
\end{lstlisting}

\subsection*{Solution 5.4}
Architecture:
\begin{enumerate}
  \item \textbf{State Machine}: LangGraph \texttt{StateGraph} with channels: \texttt{user\_id}, \texttt{verified\_identity: bool}, \texttt{active\_intent: str}, \texttt{pnr: str}, \texttt{messages: Annotated[list, add\_messages]}.
  \item \textbf{Supervisor Gating}: Triage Supervisor inspects intent. If intent requires identity verification (e.g., compensation or flight changes), the supervisor checks \texttt{verified\_identity}. If false, routing diverts to an Auth Challenge node before the Booking/Compensation workers can be invoked.
  \item \textbf{Compensation HITL Gate}: Any compensation payout exceeding \$200 triggers an \texttt{interrupt\_before=["disburse\_payment"]}, notifying a human agent via Zendesk/Slack.
\end{enumerate}

\section{Further Reading \& Citations}
\begin{itemize}[leftmargin=18pt]
  \item Yao, S., et al. (2022). \emph{ReAct: Synergizing Reasoning and Acting in Language Models}. International Conference on Learning Representations (ICLR).
  \item Wang, L., et al. (2023). \emph{Plan-and-Solve Prompting: Improving Zero-Shot Chain-of-Thought Reasoning by Large Language Models}. Annual Meeting of the Association for Computational Linguistics (ACL).
  \item Shinn, N., et al. (2023). \emph{Reflexion: Language Agents with Verbal Reinforcement Learning}. Advances in Neural Information Processing Systems (NeurIPS).
\end{itemize}
